% =====================================================================
%  Section: Related Work  (IJCAI 2027 revision: distilled lineage)
% =====================================================================
\section{Related Work}
\label{sec:related}

We organize prior work along the four axes that motivate our modular
decomposition, tracing the developmental lineage of each thread and
positioning Canvas3D at their intersection.

\subsection{From Pixel-Level Stylization to Painting Animation}

The earliest approaches to painting animation treat the input as a
\emph{pixel grid to be stylized over time}. Static image style transfer,
pioneered by feature-correlation methods~\cite{gatys2016image} and extended
by fast feed-forward networks~\cite{johnson2016perceptual,huang2017arbitrary}
and transformer-based stylizers~\cite{deng2021stytr2}, established that
artistic appearance can be re-injected into arbitrary content. Video style
transfer then addressed temporal coherence through optical-flow
regularization~\cite{ruder2016artistic,chen2017coherent,wang2020consistent},
multi-channel correlation~\cite{deng2021arbitrary}, and content deformation
fields~\cite{ouyang2024codef}. A second lineage pursues
\emph{stroke-level synthesis}---reproducing the physical painting process as
a time-lapse~\cite{zhao2020painting,zou2021stylized,liu2021neuralstrokes}---but
such methods reconstruct \emph{how a painting was made} rather than
\emph{how its subjects could move}. A third, narrower thread drives
\emph{portrait motion} on painted faces~\cite{guo2024liveportrait,tan2024edtalk}
but cannot animate full-body figures or background elements. The diffusion
era introduced inversion-based stylization~\cite{zhang2023inversion} and
video-consistent editing~\cite{ceylan2023pix2video,chai2023stablevideo,geyer2023tokenflow,liang2024flowvid},
which preserve appearance under text-driven edits but provide no geometric
handle on object pose. The recent InkSpirit~\cite{zhang2025inkspirit} and
StyleMaster~\cite{ye2025stylemaster} push appearance fidelity further yet
remain confined to 2D image translation. Across this lineage, the persistent
gap is the absence of \emph{object-level, 3D-grounded motion control}:
appearance is preserved, but spatial structure is not.

\subsection{3D Scene Representations and Single-Image Lifting}

Three-dimensional representations evolved from implicit radiance
fields~\cite{mildenhall2020nerf} to explicit primitives. 3D Gaussian
Splatting~\cite{kerbl20233dgs} represents a scene as a set of anisotropic
Gaussians with differentiable rasterization, enabling real-time
photorealistic rendering; 2D Gaussian Splatting~\cite{huang20242dgs}
collapses the third scale dimension into a surface normal, improving
geometric fidelity for surface-like content. Crucially, the field has
recently shifted from \emph{per-scene optimization} to \emph{feed-forward
single-image lifting}: generative-prior methods~\cite{tang2024dreamgaussian}
predict 3D assets from one image, while state-of-the-art feed-forward
predictors~\cite{wang2024triposplat,charatan2024pixelsplat,jiang2025anysplat}
regress Gaussian parameters directly, with two-view and unconstrained-view
variants broadening the input regime. These advances make it feasible to
lift individual painting regions into 3D, but none addresses the downstream
challenge of \emph{composing multiple such objects into a coherent,
artistically controllable scene}---exactly the gap our framework targets.

\subsection{Sketch and Line-Art as a Composition Interface}

Hand-drawn sketches have a long history as an intuitive interface for shape
prototyping~\cite{olson2011sketch}. Recent work demonstrates that
contour-level spatial reasoning is tractable with modern vision models:
ShadowDraw~\cite{luo2026shadowdraw} generates compositional shadow art from
object inputs, LineArt~\cite{wang2024lineart} performs training-free
appearance transfer onto design drawings via diffusion, and region-aware
simplification~\cite{nguyen2024region} and volumetric animation
toolkits~\cite{foxgieg2024lightning} illustrate the expressive potential of
hand-drawn control. In diffusion-based generation, sketches typically serve
as a \emph{stylistic} or \emph{boundary} condition (e.g., ControlNet-style
edge injection~\cite{zhang2023controlnet}). Canvas3D departs from this
convention by treating line-art as a \emph{spatial constraint} that
resolves the which-object-where problem in 3D scene assembly: a sketch is
not merely an appearance hint but a target silhouette that a 3D object is
geometrically registered to.

\subsection{Controllable Image-to-Video Generation}

The leap from text-to-image to text-to-video was powered by diffusion
transformers~\cite{peebles2023dit} and video diffusion
models~\cite{ho2022vdm,ho2022video}. A family of scalable
backbones~\cite{yang2025cogvideox,wan2025,ltxvideo2024,blattmann2023svd}
now supports high-fidelity image-to-video synthesis. Controllability has
been pursued through auxiliary conditioning:
ControlNet~\cite{zhang2023controlnet} injects edge, depth, and pose maps
into the diffusion U-Net, with video extensions~\cite{wang2024easycontrol}
and efficiency-oriented variants~\cite{peng2024controlnext} broadening its
applicability. Cinematic shot design~\cite{xing2025motioncanvas} and
clipart animation~\cite{wu2025aniclipart} demonstrate compelling 3D-aware
control without explicit 3D supervision. Despite these advances, existing
I2V control methods are confined to \emph{per-frame conditioning
signals}: they cannot reason about 3D object identity, depth ordering, or
occlusion across frames, and consequently drift under viewpoint change.
Canvas3D addresses this by grounding all control signals in an explicit,
object-level 3D scene that is composited, rendered, and then fed into the
video diffusion backbone.

\subsection{Painting-Specific Animation Pipelines}

A small but growing body of work targets paintings directly. Domain-adapted
detectors~\cite{chen2025pldsyolo} and art-segmentation
datasets~\cite{saleh2016metric} establish that natural-image segmentation
backbones~\cite{jocher2023yolov8,wang2024yolov9,wang2024yolov10,tian2025yolov12}
and foundation segmentation models~\cite{kirillov2023sam,ravi2024sam2}
require art-specific fine-tuning due to brushstroke texture and intentional
boundary blending. Chinese ink-wash style transfer has been explored via
instance-aware coherence~\cite{liang2021instance} and staged diffusion
guidance~\cite{liu2026chipdiff}. Most relevant to our setting,
``Printing the Underdetermined''~\cite{ming2026printing} translates
paintings to 3D representations via video generative models but operates
on the entire image jointly rather than at the object level, while
OilPaintingRevive~\cite{oilpaintingrevive2026} proposes a multi-stage
pipeline combining 2DGS with video refinement, but its skeleton-binding
network requires per-scene optimization that limits compositional
flexibility. Canvas3D advances beyond these works along four
axes---\emph{object-level} rather than image-level manipulation,
\emph{line-art} as a compositional control interface,
\emph{depth-aware inpainting} for disocclusion, and
\emph{multi-modal geometric conditioning} for video synthesis---yielding
stronger 3D consistency and explicit compositional control.
